In the active case, the volume \(V_{int}\) is determined not only by diffusion but also by advection. A deviation \(\delta\theta\) creates different separation between particles in the \(\hat n \parallel \hat p\) direction and the \(\hat n \perp \hat p\) direction. This gives us
\begin{equation}
    \hat P = \ins{\cos\theta,\sin\theta} \approx \ins{1-\frac{1}{2}\delta\theta^2,\delta\theta}
\end{equation}
and
\begin{align}
    \delta x_\perp \sim v_0\tau\sin\delta\theta \sim \tau\delta\theta && \delta x_\parallel \sim v_0\tau\delta\theta^2 \sim \tau\delta\theta^2
\end{align}
and the volume
\begin{equation}
    V_e \sim \omega_\perp^{d-1}\omega\parallel
\end{equation}
From a combination of advection and diffusion:
\begin{equation}
    \omega_\perp \sim \delta x_\perp + D_\perp\tau^{1/2} \sim \tau\delta\theta + D_\perp\tau^{1/2} \sim \tau^{\gamma_\perp}
\end{equation}
and
\begin{equation}
    \omega_\parallel \sim \delta x_\parallel + D_\parallel\tau^{1/2} \sim \tau\delta\theta^2+D_\parallel\tau^{1/2}\sim\tau^\gamma
\end{equation}
if we define \(\delta\theta\sim\tau^\gamma\) then we also find
\begin{equation}
    \delta\theta \sim \sqrt{\frac{\tau}{V_e}} \sim \frac{\tau^{1/2}}{\sqrt{\omega_\perp^{d-1}\omega\parallel}}\sim\tau^\gamma
\end{equation}
We thus have to solve three equations for the gammas:
\begin{align}
    2\gamma &= 1-\gamma_\parallel-\ins{d-1}\gamma_\perp \\
    \gamma_\perp &= \max\ins{1+\gamma,\frac{1}{2}} \\
    \gamma_\parallel &= \max\ins{1+2\gamma,\frac{1}{2}}
\end{align}
we find \(\gamma < 0\) for all \(d>1\). Specifically for \(d=2\) we find
\begin{align}
    \gamma = -\frac{1}{5} && \gamma_\perp = \frac{4}{5} && \gamma_\parallel = \frac{3}{5}
\end{align}
This means the order is stable, thanks to the `quick smearing' in the direction perpendicular to \(\vec{P}\). A couple of notes:
\begin{itemize}
    \item The stability of polar phase is still being researched in a more rigorous manner.
    \item Current research shows that active polar order has additional properties. For example, a boundary effect can affect order in bulk (Yariv Kafri).
\end{itemize}
\subsection{Waves in Active Polar Matter}
Remembering the equations
\begin{align}\label{polarwaves1}
    \diffp{P}{t} = -\partial_\alpha j_\alpha && j_\alpha = -\gamma_\rho\partial_\alpha\frac{\delta F}{\delta\rho}+v_0\rho P_\alpha && \diffp{P_\alpha}{t} + \lambda\ins{P_\beta\partial_\beta} = -\gamma_P \frac{\delta F}{\delta P_\alpha}
\end{align}
We'll examine small deviations around the homogenous stable state
\begin{align}
    \rho = \rho_0 + \delta\rho(t,\vec{r}) && \vec{P} = \vec{P_0} + \delta\vec{P}(t,\vec{r})
\end{align}
for
\begin{equation}
    \frac{\delta \rho}{\rho_0},\frac{\lvert\vec{P}\rvert}{\lvert\vec{P}\rvert} \ll 1
\end{equation}
we'll linearize the equations and keep leading elements in the derivatives (first order). For simplicity we'll assume \(\vec{P} = \hat{x} + \delta P\hat{y}\implies\) the steady state is completely ordered and only rotations of the polarizations are possible. In this case, there is a steady current \(\vec{j_0}=\rho_0v_0\hat{x}\). In this case we won't interest ourselves with the Landau elements of the free energy (\(\frac{a}{2}P^2+\frac{b}{4}P^4\)).
\begin{align}
    F = \int\dd^2 r \left[\frac{\chi}{2}\ins{\frac{\delta\rho}{\rho_0}}^2+\frac{1}{2}k\partial_\alpha P_\beta\partial_\alpha P_\beta + \ins{\omega_0-\omega_1\frac{\delta\rho}{\rho_0}}\partial_\alpha P_\alpha\right]
\end{align}
What does the last element mean? From integration by parts
\begin{align}
    F &= \int\dd^2 r \left[\dots - P_\alpha\partial_\alpha\ins{\omega_0-\omega_1\frac{\delta\rho}{\rho_0}}\right] + \text{boundary element}\\
    &= \int\dd^2 r \left[\dots + \frac{\omega_1}{\rho}\vec{P}\cdot\nabla\rho\right]
\end{align}
\(\omega_0\) only contributes a boundary element. \(\omega_1\) determines whether \(\vec{P}\) likes to order itself with or against the gradients in the density.
% figure
We'll develop the equation in \ref{polarwaves1} to first order in \(\delta\rho, \delta P\) and up to derivatives of first order and find
\begin{align}
    \partial_t\delta\rho = -\partial_\alpha\ins{v_0\delta\rho P_\alpha^0 + v_0\rho_0\delta P_\alpha}
\end{align}
where
\begin{align}
    \ins{\partial_t+v_0\partial_x}\delta\rho \approx -v_0\rho_0\partial_y\delta P \\
    \ins{\partial_t + \lambda\partial_x}\delta P \approx -\gamma_P\omega_1\partial_y\frac{\delta\rho}{\rho_0}
\end{align}
because
\begin{equation}
    \frac{\delta F}{\delta P_\alpha} = -k\partial_{\beta\beta}\delta P_\alpha + \frac{\omega_1}{\rho_0}\partial_\alpha\delta\rho
\end{equation}
If the system is homogenous in the \(y\) direction, we find that \(\delta P, \delta\rho\) are independent up to first order. One moves at velocity \(v_0\) and the other at velocity \(\lambda\). If the system is homogenous in the \(x\) direction, we find
\begin{align}
    \partial_t\delta\rho = -v_0\rho_0\partial_y\delta P \\
    \partial_t\delta P = -\gamma_P\omega_1\partial_y\frac{\delta\rho}{\rho_0}
\end{align}
which gives us
\begin{equation}
    \partial_{tt}\delta\rho = v_0\gamma_P\omega_1\partial_{yy}\delta\rho 
\end{equation}
for \(\omega_1 > 0\) we found a wave equation with the velocity 
\begin{equation}
    c = \sqrt{\omega_1\gamma_P v_0}
\end{equation}
Changes in the density reorganize \(\vec{P}\) and changes in \(\vec{P}\) transport the density because of the activity. Note: the higher order terms in derivatives will lead to decay of the waves, i.e to diffusion.
\newpage
\section{Nematic Active Matter}
\subsection{The Nematic Order Parameter}
So far we dealt with polar order \(\vec{P} = \langle\hat l\rangle\) with a distinct direction. Sometimes there is only a distinct axis. This kind of order is called nematic. The order parameter is the Nematic Tensor:
\begin{equation}
    Q_{\alpha\beta} = \left\langle\hat{l_\alpha}\hat{l_\beta}-\frac{\delta_{\alpha\beta}}{d}\right\rangle = Q\ins{n_\alpha n_\beta-\frac{\delta_{\alpha\beta}}{d}}
\end{equation}
Its properties are
\begin{enumerate}
    \item Invariant to the substitution \(\hat{n} \iff -\hat{n}\)
    \item It is symmetric and traceless
\end{enumerate}
For \(d=2\) and complete order along the \(\hat{x}\) axis:
\begin{equation}
    Q_{\alpha\beta} = \begin{pmatrix}
        \frac{1}{2} && 0 \\ 0 && -\frac{1}{2}
    \end{pmatrix}
\end{equation}
in general, in the direction \(\theta \rightarrow \hat{n}=\ins{\cos\theta,\sin\theta}\)
\begin{equation}
    Q_{\alpha\beta} = \frac{1}{2}\begin{pmatrix}
        \cos\ins{2\theta} && \sin\ins{2\theta} \\ \sin\ins{2\theta} && -\cos\ins{2\theta}
    \end{pmatrix}
\end{equation}
with the angle being \(2\theta\) representing the fact that \(\hat{n}\iff-\hat{n}\) (\(\theta\iff\theta+\pi\)). A physical theory of nematic matter can be also written in terms of \(\hat{n}\) while maintaining the symmetry for \(\hat{n}\iff-\hat{n}\) and a scalar parameter \(Q\). In practice, it is customary to sometimes use \(\vec{P}\) instead of \(\hat{n}\). \underline{Remember} this does not mean there is a preferred direction. Even if we can write \(\left\langle\hat{l}_\alpha\right\rangle\) we can write
\begin{equation}
    Q_{\alpha\beta} = P_\alpha P_\beta -P_\mu P_\mu \frac{\delta_{\alpha\beta}}{d}
\end{equation}
and in nematic theory the symmetry \(\vec{P}\iff-\vec{P}\) must be maintained. What does that mean? Let's return to the equations.
\begin{equation}
    F = \int \dd \vec{r}\left[\frac{a}{2}P^2+\frac{b}{4}P^4+\frac{K}{2}\partial_\alpha P_\beta\partial_\alpha P_\beta-\cancel{\omega(\rho)\partial_\alpha P_\alpha} + \dots\right]
\end{equation}
The final element does not fit the symmetry we require so we cannot use it in nematic theory.
\begin{align}
    \diffp{\rho}{t} = -\partial_\alpha j_\alpha && j_\alpha = -\gamma_\rho\partial_\alpha\frac{\delta F}{\delta\rho} + \cancel{v_0\rho P_\alpha}
\end{align}
the final element in the right hand equation also cannot be used for the same reason. Finally
\begin{equation}
    \diffp{P_\alpha}{t} + \cancel{\lambda \ins{P_\mu\partial_y}P_\alpha} = -\gamma_P\rho \frac{\delta F}{\delta P_\alpha}
\end{equation}
The unintuitively cancelled element is \(\ins{\vec{P}\cdot\nabla}\vec{P}\). Why is it unallowed? Intuitively it comes from advection which no longer exists. Mathematically, we'll multiply by \(P_\beta\) and find
\begin{align}
    P_\beta\partial_t P_\alpha + \lambda P_\beta\ins{P_\mu \partial_\mu}P_\alpha = \dots \\
    \partial_t Q_{\alpha\beta} + \cancel{\lambda\ins{P_\mu\partial_\mu}Q_{\alpha\beta}} = \dots
\end{align}
What is the active contribution? We need to take into consideration active currents. How can we build it from \(Q_{\alpha\beta}\)? We need to apply a vector on it. Without a defined direction, the only vector is \(\vec{\nabla}\):
\begin{equation}
    \vec{j} = -\gamma_\rho\vec{\nabla}\frac{\delta F}{\delta\rho}-\Lambda\vec{\nabla}\cdot\vec{Q}
\end{equation}